\section*{Introduction: The Clash of Titans - Einstein's Realism vs. Quantum Indeterminacy}

At the heart of 20th-century physics lies a deep philosophical rift, a schism that pitted two seemingly irreconcilable worldviews against each other.
On one side, Albert Einstein, whose physical intuition was deeply rooted in classical determinism and the philosophy of Spinoza, held the belief that 
the universe is, at its most fundamental level, knowable and orderly. For Einstein, the physical properties of an object, such as its position or momentum, 
existed objectively, with definite values, regardless of whether an observer measured them. This worldview directly clashed with the principles of the emerging 
quantum theory, particularly the Copenhagen Interpretation, which embraced an intrinsic indeterminacy and a fundamental randomness as inherent features of nature.

In 1935, Einstein, along with his collaborators Boris Podolsky and Nathan Rosen, launched his most formidable challenge to the new physics in a seminal paper. 
The argument, now known as the EPR paradox, was not a mere technicality but a frontal assault on the conceptual foundations of quantum mechanics. It is crucial 
to understand that the EPR paradox was not an attempt to refute the predictive power of quantum mechanics, but rather to question its claim to be a complete theory. 
The very title of the paper, ``Can Quantum-Mechanical Description of Physical Reality Be Considered Complete?'', frames the argument in precisely these terms. Einstein
accepted that the theory's calculations were correct, but he argued that its probabilistic nature was not a fundamental feature of reality, but rather a reflection of 
our ignorance about a deeper level of description.

\subsection*{The Criterion of Reality}

To formalize their attack, EPR introduced an explicit and, on its face, irrefutable criterion for defining an element of physical reality. The EPR Criterion of Reality states:

\textit{if, without in any way disturbing a system, we can predict with certainty (i.e., with probability equal to unity) the value of a physical quantity, then there 
exists an element of physical reality corresponding to this physical quantity.}

To illustrate how quantum mechanics broke this criterion of Reality, the paper showed how the momentum and position operators could not be determined simultaneously.
Given a state:
$$ \psi = e^{\frac{2\pi i}{h}p_{0}x} $$
by applying the momentum operator:
$$ p = (h/2\pi i)\partial/\partial x $$
we obtain
$$ p\psi = (h/2\pi i)\partial\psi/\partial x = p_{0}\psi. $$
Thus, the value of the momentum is known with certainty to be $p_0$. However, for the position, it is different, since, where a is a constant, there is no specific position:
$$ q\psi = x\psi \ne a\psi $$
One can only speak of the probability of finding the particle between a and b:
$$ P(a,b) = \int_{a}^{b}\overline{\psi}\psi dx = \int_{a}^{b}dx = b-a. $$
The usual conclusion from this in quantum mechanics is that when the momentum of a particle is known, its coordinate has no physical reality. From this it follows that either
(1) the quantum-mechanical description of reality given by the wave function is not complete or (2) when the operators corresponding to two physical quantities do not commute, 
the two quantities cannot have simultaneous reality.

\subsection*{The Principle of Locality}

The principle of locality is one of the most intuitive and fundamental ideas in physics. It states that an object can only be directly influenced by its immediate surroundings. 
Any effect or influence at a distance cannot be instantaneous; it must propagate through space at a finite speed, which cannot exceed the speed of light in a vacuum, c. For Einstein,
locality was a sacred axiom; its violation would imply the disintegration of the causal structure of the universe, where cause precedes effect in a well-defined spacetime order, making the
formulation of consistent physical laws impossible.

\subsection*{The EPR Experiment}

An experiment can be performed to better explain this situation. The original 1935 thought experiment involves a system of two particles that interact briefly and then separate in opposite
directions. Quantum mechanics allows this system to be prepared in an entangled state such that, due to conservation laws, its total momentum is known (e.g., zero) and its relative distance
is also known at all times. Two observers, conventionally named Alice and Bob, each receive one of the particles and can perform measurements on them at a great distance.

The argument unfolds in two alternative steps:
\begin{enumerate}
    \item \textbf{Momentum Measurement:} Alice decides to measure the momentum of her particle (particle 1). Due to the law of conservation of momentum, the instant she obtains a result, 
    she can predict with absolute certainty the momentum of Bob's particle (particle 2). Applying the Criterion of Reality, and assuming locality (Alice's measurement cannot disturb Bob's 
    particle), it is concluded that the momentum of particle 2 is an ``element of reality''.

    \item \textbf{Position Measurement:} Alternatively, Alice could have chosen to measure the position of her particle. Since their relative position is known, this measurement would allow 
    her to predict with certainty the position of Bob's particle. Again, invoking the Criterion of Reality and locality, it is concluded that the position of particle 2 must also be an 
    ``element of reality''.
\end{enumerate}

The strength of the argument lies in the observer's freedom of choice. Alice is free to decide which quantity to measure. Since her choice, made at her location, cannot instantaneously 
influence the physical reality of Bob's distant particle (principle of locality), both properties—the position and momentum of Bob's particle—must have been elements of reality simultaneously, 
from the moment the particles separated. Alice's measurement would simply reveal one of these pre-existing values.

Here the central contradiction arises. The conclusion that Bob's particle simultaneously possesses a definite value for its position and its momentum is in direct conflict with one of the
pillars of quantum mechanics: Heisenberg's Uncertainty Principle, which forbids two incompatible (or non-commuting) observables, such as position and momentum, from having well-defined values 
at the same time. Faced with this dilemma, EPR reached an inevitable conclusion from their perspective: if local realism is correct, then the particles must possess properties that quantum 
mechanics is unable to describe in its formalism. Therefore, the description of reality provided by the wave function must be incomplete.

\section*{Bohr's Direct Rebuttal to the EPR Criterion of Reality}

Niels Bohr's response, published in the same journal a few months later, does not attempt to find a flaw in EPR's logical chain. To understand Bohr's response, it is imperative to first
understand his central concept: complementarity. Complementarity postulates that certain pairs of classical concepts (like particle/wave, or position/momentum) are necessary for a complete 
description of quantum phenomena, but are mutually exclusive in any single experiment. They are not contradictory properties of an underlying object, but complementary aspects of a phenomenon
that can only be revealed through mutually incompatible experimental arrangements.

For Bohr, a quantum phenomenon is an indivisible whole that includes both the object and the specific measuring apparatus used to investigate it. The experimental context is not a passive 
backdrop, but an active and constitutive part of the phenomenon. Measuring position creates a ``position phenomenon''; measuring momentum creates a ``momentum phenomenon''. They are not two 
different views of the same underlying reality, but two different and mutually exclusive phenomenal realities.

\section*{Schrödinger's 1935 Response to the EPR Paradox}

A few months after the publication of the EPR paper, Erwin Schrödinger, one of the founding architects of quantum mechanics, published his own response. He argues that the situation described by
EPR is not a simple puzzle, but rather ``the characteristic trait of quantum mechanics, the one that enforces its entire departure from classical lines of thought''. To describe this situation,
he coined the term \textit{Verschränkung}, which translates to ``entanglement''.

To rigorously demonstrate this phenomenon, he used a mathematical formalism now known as the Schmidt decomposition. He showed that for any entangled state of two systems, $\Psi(x,y)$,
one can write an expansion of the form $\Psi(x,y)=\Sigma_n c_n g_n(x)f_n(y)$.
